% ============================================================
%  Experiment 5 --- Interfacing LCD to Altera DE2 Board
% ============================================================
\experimentheader{5}{LCD Interfacing}{Interfacing LCD to Altera DE2 Board}

\begin{objectives}
\begin{enumerate}
  \item To use Verilog code as a design entry method in Altera Quartus II Software.
  \item To interface an LCD module to the Altera DE2 board.
\end{enumerate}
\end{objectives}

\begin{equipment}
\begin{itemize}
  \item Altera Quartus II software
  \item Altera DE2 board
\end{itemize}
\end{equipment}

\begin{introduction}
In recent years, LCDs have largely replaced simple LED displays: they can show numbers, characters, and graphics, and are relatively easy to program. Table~\ref{tab:exp5-pins-lcd} lists the LCD pins relevant to programming.

The \textit{enable} (E) pin is used by the LCD to latch the information presented on its data pins: a high-to-low transition on E must be applied for the LCD to latch the data, and this pulse must be at least 450\,ns wide. The required timing is shown in Figure~\ref{fig:exp5-timing}.

Two important internal registers are selected via the RS pin. When $\text{RS}=0$, the \textit{instruction/command} register is selected, allowing commands to be sent to the LCD (Table~\ref{tab:exp5-commands} lists the available command codes). When $\text{RS}=1$, the \textit{data} register is selected, allowing ASCII character data to be written for display, sent through data pins D0--D7.

It is good practice to check the busy flag before writing to the LCD: with $\text{RS}=0$ and $\text{R}/\overline{\text{W}}=1$, the busy flag appears on D7. While D7\,=\,1 the LCD is busy with internal operations and will not accept new data; once D7\,=\,0, the LCD is ready to receive the next byte.
\end{introduction}

\begin{boxedtable}
\centering
\begin{tabular}{clcl}
\toprule
Pin & Symbol & I/O & Description \\
\midrule
1  & $V_{ss}$ & --- & Ground \\
2  & $V_{cc}$ & --- & $+5\,\text{V}$ power supply \\
3  & $V_{ee}$ & --- & Power supply to control contrast \\
4  & RS   & I   & $\text{RS}=0$: command register; $\text{RS}=1$: data register \\
5  & R/$\overline{\text{W}}$ & I & $0$: write; $1$: read \\
6  & E    & I/O & Enable (latches data on high-to-low transition) \\
7  & DB0  & I/O & \multirow{8}{*}{8-bit data bus} \\
8  & DB1  & I/O & \\
9  & DB2  & I/O & \\
10 & DB3  & I/O & \\
11 & DB4  & I/O & \\
12 & DB5  & I/O & \\
13 & DB6  & I/O & \\
14 & DB7  & I/O & \\
\bottomrule
\end{tabular}
\captionof{table}{Pin description for the LCD module.}
\label{tab:exp5-pins-lcd}
\end{boxedtable}

\begin{boxedfigure}
\centering
\resizebox{0.98\linewidth}{!}{%
\begin{tikzpicture}[
  x=1.05cm, y=0.9cm,
  sig/.style={very thick, draw=inkblack},
  lbl/.style={font=\small\plexmedium, anchor=east},
]
  % row y-positions
  \def\yD{6}\def\yE{4}\def\yRW{2}\def\yRS{0}

  \node[lbl] at (-0.3,\yD)  {D0--D7};
  \node[lbl] at (-0.3,\yE)  {E};
  \node[lbl] at (-0.3,\yRW) {R/$\overline{\text{W}}$};
  \node[lbl] at (-0.3,\yRS) {RS};

  % D0-D7 line: low, then hexagon "Data" pulse, then low
  \draw[sig] (0,\yD) -- (3.4,\yD);
  \draw[sig] (3.4,\yD) -- (3.8,\yD+0.4) -- (7.6,\yD+0.4) -- (8.0,\yD) -- (7.6,\yD-0.4) -- (3.8,\yD-0.4) -- cycle;
  \node[font=\small] at (5.7,\yD) {Data};
  \draw[sig] (8.0,\yD) -- (11,\yD);
  % tD arrow
  \draw[<->, thin] (3.4,\yD-0.9) -- (4.6,\yD-0.9);
  \node[font=\scriptsize] at (4.0,\yD-1.25) {$t_D$};
  \draw[dotted] (3.4,\yD-1.4) -- (3.4,\yE+0.55);
  \draw[dotted] (4.6,\yD-1.4) -- (4.6,\yE+0.55);

  % E line: low, high pulse, low
  \draw[sig] (0,\yE) -- (3.4,\yE) -- (3.4,\yE+0.55) -- (7.6,\yE+0.55) -- (7.6,\yE) -- (11,\yE);
  \draw[dotted] (7.6,\yE+0.9) -- (7.6,\yRW-0.55);

  % R/W line: low pulse near start (setup), then returns; tAS / tAH markers
  \draw[sig] (0,\yRW) -- (2.9,\yRW) -- (2.9,\yRW+0.5) -- (3.9,\yRW+0.5) -- (3.9,\yRW) -- (7.1,\yRW) -- (7.1,\yRW+0.5) -- (8.1,\yRW+0.5) -- (8.1,\yRW) -- (11,\yRW);
  \draw[<->, thin] (2.9,\yRW-0.35) -- (3.4,\yRW-0.35);
  \node[font=\scriptsize] at (3.15,\yRW-0.62) {$t_{AS}$};
  \draw[<->, thin] (7.6,\yRW-0.35) -- (8.1,\yRW-0.35);
  \node[font=\scriptsize] at (7.85,\yRW-0.62) {$t_{AH}$};

  % RS line: low, then step high covering the enable pulse, then low
  \draw[sig] (0,\yRS) -- (3.4,\yRS) -- (3.4,\yRS+0.55) -- (7.6,\yRS+0.55) -- (7.6,\yRS) -- (11,\yRS);
\end{tikzpicture}
}

\vspace{4pt}
\footnotesize
$t_D$ = data output delay time \\
$t_{AS}$ = setup time prior to E (going high) for RS and R/$\overline{\text{W}}$, minimum 140\,ns \\
$t_{AH}$ = hold time after E falls, for RS and R/$\overline{\text{W}}$, minimum 10\,ns \\
\textit{Note: a read operation requires an L-to-H pulse on the E pin.}
\normalsize
\captionof{figure}{LCD write-cycle timing diagram.}
\label{fig:exp5-timing}
\end{boxedfigure}

\begin{boxedtable}
\centering
\begin{tabular}{cl}
\toprule
Code (hex) & Command to LCD instruction register \\
\midrule
1  & Clear display screen \\
2  & Return home \\
4  & Decrement cursor (shift cursor left) \\
6  & Increment cursor (shift cursor right) \\
5  & Shift display right \\
7  & Shift display left \\
8  & Display off, cursor off \\
A  & Display off, cursor on \\
C  & Display on, cursor off \\
E  & Display on, cursor blinking \\
F  & Display on, cursor blinking \\
10 & Shift cursor position left \\
14 & Shift cursor position right \\
18 & Shift entire display left \\
1C & Shift entire display right \\
80 & Force cursor to beginning of line 1 \\
C0 & Force cursor to beginning of line 2 \\
38 & 2 lines, $5\times7$ dot matrix \\
\bottomrule
\end{tabular}
\captionof{table}{LCD command codes.}
\label{tab:exp5-commands}
\end{boxedtable}

\begin{procedure}

\textbf{Activity 1 --- Displaying lookup-table data on the LCD}\vspace{2pt}

\textit{Objective: display fixed text, stored in a lookup table, on the LCD.}

\begin{enumerate}
  \item Open Quartus II.
  \item Write the Verilog modules given in Listing~\ref{lst:exp5-lcdt} (\texttt{lcd\_t}) and Listing~\ref{lst:exp5-lcdctrl} (\texttt{LCD\_Controller}) to send commands and data to the LCD.
  \item \texttt{lcd\_t} uses a lookup table (LUT) to store the sequence of commands and characters to send to the LCD.
  \item Create a reusable symbol for each module, as in Experiment 3.
  \item Draw the schematic in Figure~\ref{fig:exp5-schematic}, connecting \texttt{lcd\_t}'s outputs to \texttt{LCD\_Controller}'s inputs, and \texttt{LCD\_Controller}'s outputs to the LCD interface pins.
  \item Assign pin numbers as shown in Table~\ref{tab:exp5-pinassign}.
  \item Download the program to the Altera DE2 board.
  \item Discuss the result based on your observation of the LCD.
\end{enumerate}

\begin{boxedfigure}
\centering
\resizebox{0.98\linewidth}{!}{%
\begin{tikzpicture}[node distance=14mm and 18mm]
  \node[io] (clk) {CLK};
  \node[io, below=6mm of clk] (rst) {RST};
  \node[block, right=of clk, minimum width=3cm, minimum height=2.6cm, yshift=-8mm] (lcdt) {lcd\_t\\[2pt]\scriptsize CLK, RST $\to$\\ mLCD\_DATA,\\ mLCD\_RS, mLCD\_Start};
  \node[block, right=42mm of lcdt, minimum width=3.2cm, minimum height=2.8cm] (lcdc) {LCD\_Controller\\[2pt]\scriptsize iDATA, iRS, iStart $\to$\\ LCD\_DATA, LCD\_RW,\\ LCD\_EN, LCD\_RS, LCD\_POW};
  \node[io, right=of lcdc, yshift=10mm] (lcddata) {LCD\_DATA[7:0]};
  \node[io, right=of lcdc, yshift=2mm] (lcden) {LCD\_EN};
  \node[io, right=of lcdc, yshift=-6mm] (lcdrs) {LCD\_RS};
  \node[io, right=of lcdc, yshift=-14mm] (lcdpow) {LCD\_POW};

  \draw[arr] (clk) -- (lcdt);
  \draw[arr] (rst) -- (lcdt);
  \draw[arr] (lcdt) -- (lcdc) node[midway, above, font=\tiny] {mLCD\_DATA, mLCD\_RS, mLCD\_Start};
  \draw[arr] (lcdc.east) ++(0,10mm) -- (lcddata);
  \draw[arr] (lcdc.east) ++(0,2mm) -- (lcden);
  \draw[arr] (lcdc.east) ++(0,-6mm) -- (lcdrs);
  \draw[arr] (lcdc.east) ++(0,-14mm) -- (lcdpow);
\end{tikzpicture}
}
\captionof{figure}{Schematic diagram for the LCD interface: \texttt{lcd\_t} (lookup-table sequencer) feeding \texttt{LCD\_Controller} (bus-timing controller).}
\label{fig:exp5-schematic}
\end{boxedfigure}

\begin{boxedtable}
\centering
\begin{tabular}{clcl}
\toprule
\# & Node name & Direction & Location \\
\midrule
1  & LCD\_DATA[7] & Output & PIN\_H3 \\
2  & LCD\_DATA[6] & Output & PIN\_H4 \\
3  & LCD\_DATA[5] & Output & PIN\_J3 \\
4  & LCD\_DATA[4] & Output & PIN\_J4 \\
5  & LCD\_DATA[3] & Output & PIN\_H2 \\
6  & LCD\_DATA[2] & Output & PIN\_H1 \\
7  & LCD\_DATA[1] & Output & PIN\_J2 \\
8  & LCD\_DATA[0] & Output & PIN\_J1 \\
9  & LCD\_EN      & Output & PIN\_K3 \\
10 & LCD\_POW     & Output & PIN\_L4 \\
11 & LCD\_RS      & Output & PIN\_K1 \\
12 & LCD\_RW      & Output & PIN\_K4 \\
13 & CLK          & Input  & PIN\_N2 \\
14 & RST          & Input  & PIN\_G26 \\
\bottomrule
\end{tabular}
\captionof{table}{Pin assignment for the LCD interface.}
\label{tab:exp5-pinassign}
\end{boxedtable}

\begin{codepanel}{lcd\_t.v --- Appendix 1: lookup-table sequencer}
\begin{minted}{verilog}
module lcd_t (
    // Host side
    CLK, RST, mLCD_Done,
    mLCD_DATA, mLCD_RS, mLCD_Start
);
  // Host side
  input  CLK, RST, mLCD_Done;
  // Internal wires/registers
  output [7:0] mLCD_DATA;
  output mLCD_RS, mLCD_Start;
  reg  [5:0] LUT_INDEX;
  reg  [8:0] LUT_DATA;
  reg  [5:0] mLCD_ST;
  reg  [17:0] mDLY;
  reg  mLCD_Start;
  reg  [7:0] mLCD_DATA;
  reg  mLCD_RS;
  wire mLCD_Done;

  parameter LCD_INITIAL = 0;
  parameter LCD_LINE1   = 5;
  parameter LCD_CH_LINE = LCD_LINE1 + 16;
  parameter LCD_LINE2   = LCD_LINE1 + 16 + 1;
  parameter LUT_SIZE    = LCD_LINE1 + 32 + 1;

  // Sequencing state machine: send one LUT entry, wait for
  // LCD_Controller to finish (mLCD_Done), pause, then advance.
  always @(posedge CLK or negedge RST)
  begin
    if (!RST)
    begin
      LUT_INDEX  <= 0;
      mLCD_ST    <= 0;
      mDLY       <= 0;
      mLCD_Start <= 0;
      mLCD_DATA  <= 0;
      mLCD_RS    <= 0;
    end
    else if (LUT_INDEX < LUT_SIZE)
    begin
      case (mLCD_ST)
        0: begin
          mLCD_DATA  <= LUT_DATA[7:0];
          mLCD_RS    <= LUT_DATA[8];
          mLCD_Start <= 1;
          mLCD_ST    <= 1;
        end
        1: if (mLCD_Done) begin
          mLCD_Start <= 0;
          mLCD_ST    <= 2;
        end
        2: if (mDLY < 18'h3FFFE)
             mDLY <= mDLY + 1;
           else begin
             mDLY    <= 0;
             mLCD_ST <= 3;
           end
        3: begin
          LUT_INDEX <= LUT_INDEX + 1;
          mLCD_ST   <= 0;
        end
      endcase
    end
  end

  // Lookup table: LCD_INITIAL entries are setup commands
  // (function set, display on, clear, entry mode, home);
  // LCD_LINE1 / LCD_LINE2 entries are ASCII characters for
  // each line, with the MSB (bit 8) marking RS = 1 (data).
  always
  begin
    case (LUT_INDEX)
      LCD_INITIAL+0: LUT_DATA <= 9'h038;  // function set
      LCD_INITIAL+1: LUT_DATA <= 9'h00C;  // display on
      LCD_INITIAL+2: LUT_DATA <= 9'h001;  // clear display
      LCD_INITIAL+3: LUT_DATA <= 9'h006;  // entry mode set
      LCD_INITIAL+4: LUT_DATA <= 9'h080;  // cursor home, line 1
      // Line 1: "Welcome to the "
      LCD_LINE1+0:  LUT_DATA <= 9'h120;
      LCD_LINE1+1:  LUT_DATA <= 9'h157;
      LCD_LINE1+2:  LUT_DATA <= 9'h165;
      // ... remaining characters follow the same pattern ...
      LCD_LINE1+15: LUT_DATA <= 9'h120;
      LCD_CH_LINE:  LUT_DATA <= 9'h0C0;   // move to line 2
      // Line 2: "    EEM 441 Lab "
      LCD_LINE2+0:  LUT_DATA <= 9'h120;
      LCD_LINE2+4:  LUT_DATA <= 9'h145;
      // ... remaining characters follow the same pattern ...
      LCD_LINE2+15: LUT_DATA <= 9'h120;
      default:      LUT_DATA <= 9'h000;
    endcase
  end
endmodule
\end{minted}
\end{codepanel}
\vspace{-4pt}
\begin{center}\small\color{slategray}Listing \refstepcounter{equation}\theequation\label{lst:exp5-lcdt}: \texttt{lcd\_t.v} --- lookup-table LCD message sequencer.\end{center}

\begin{codepanel}{LCD\_Controller.v --- Appendix 2: bus-timing controller}
\begin{minted}{verilog}
module LCD_Controller (
    // Host side
    iDATA, iRS, iStart, oDone, CLK, RST,
    // LCD interface
    LCD_DATA, LCD_RW, LCD_EN, LCD_RS, LCD_POW
);
  parameter CLK_Divide = 16;

  // Host side
  input  [7:0] iDATA;
  input  iRS, iStart;
  input  CLK, RST;
  output reg oDone;

  // LCD interface
  output [7:0] LCD_DATA;
  output reg LCD_EN;
  output LCD_RW;
  output LCD_RS;
  output LCD_POW;

  // Internal registers
  reg [4:0] Cont;
  reg [1:0] ST;
  reg preStart, mStart;

  // Write-only interface: bypass iRS straight to LCD_RS
  assign LCD_DATA = iDATA;
  assign LCD_RW   = 1'b0;
  assign LCD_RS   = iRS;
  assign LCD_POW  = 1'b1;

  always @(posedge CLK or negedge RST)
  begin
    if (!RST)
    begin
      oDone    <= 1'b0;
      LCD_EN   <= 1'b0;
      preStart <= 1'b0;
      mStart   <= 1'b0;
      Cont     <= 0;
      ST       <= 0;
    end
    else
    begin
      // Rising-edge detect on iStart
      preStart <= iStart;
      if ({preStart, iStart} == 2'b01)
      begin
        mStart <= 1'b1;
        oDone  <= 1'b0;
      end

      if (mStart)
      begin
        case (ST)
          0: ST <= 1;                 // wait setup
          1: begin
            LCD_EN <= 1'b1;           // raise enable
            ST     <= 2;
          end
          2: if (Cont < CLK_Divide)   // hold enable pulse width
               Cont <= Cont + 1;
             else
               ST <= 3;
          3: begin
            LCD_EN <= 1'b0;           // drop enable: latches data
            mStart <= 1'b0;
            oDone  <= 1'b1;
            Cont   <= 0;
            ST     <= 0;
          end
        endcase
      end
    end
  end
endmodule
\end{minted}
\end{codepanel}
\vspace{-4pt}
\begin{center}\small\color{slategray}Listing \refstepcounter{equation}\theequation\label{lst:exp5-lcdctrl}: \texttt{LCD\_Controller.v} --- generates the LCD enable-pulse timing.\end{center}

\textbf{Activity 2 --- Displaying DIP switch input on the LCD}\vspace{2pt}

\textit{Objective: display the live 8-bit DIP switch configuration on the LCD.}

\begin{enumerate}
  \item Write Verilog code that reads the 8-bit DIP switch configuration and displays it on the LCD (e.g.\ if switch 1 is turned on, the LCD should display \texttt{1} in the corresponding position).
  \item Print and submit the Verilog code and simulation result.
\end{enumerate}

\end{procedure}
